\documentclass[11pt,a4paper]{article}
\usepackage[T1]{fontenc}
\usepackage{amsmath,amssymb,geometry}
\geometry{margin=29mm}
\title{EPPA--III: a direct profile proof of the equal-clique case}
\author{Independent optional validation; original manuscript unchanged}
\date{9 October 2026}
\begin{document}
\maketitle

\paragraph{Theorem (stronger than Lemma 3.7 for
non-homogeneous witnesses).}
Let \(G=sK_t\) be induced in a finite
\emph{non-homogeneous vertex-transitive}
EPPA witness \(H\), with \(|H|=2st\).
Then \(H\) is a generalised clique cover
of \(G\), or \(s=t=2\) and \(H\) is
the Wagner graph. No imprimitivity or
a priori decomposition of \(H\setminus G\)
is assumed.

\paragraph{Proof.}
Put \(B=H\setminus G\), \(n=st\), and
let \(C_1,\ldots,C_s\) be the components
of \(G\). Regularity and the equality of
the two part sizes imply
\[
             |E(H[B])|=|E(G)|=s\binom t2.
             \tag{1}
\]
For \(s=1\), the outside graph is \(K_t\) by
(1), and the \(S_t\)-orbits on outside
neighbourhoods have sizes \(\binom tk\le t\).
They are therefore empty, full, singleton,
or co-singleton; vertex regularity excludes
mixing orbits, giving the four diagonal
connection types. For \(t=1\), equation (1)
makes \(B\) independent, and the \(S_s\)-orbits
on outside neighbourhoods have sizes
\(\binom sr\le s\). Regularity again excludes
mixing degrees, giving empty, complete,
matching, or co-matching cross-connections.
These are all generalised clique covers.
Henceforth assume \(s,t\ge2\).
For \(w\in B\) record its neighbourhood
profile \(P(w)=(N_H(w)\cap C_i)_{i=1}^s\).
Every permutation in
\(W=\operatorname{Aut}(G)=S_t\wr S_s\)
extends to an automorphism of \(H\)
preserving \(G\) and \(B\). Therefore
\[
                      |W P(w)|\le st.
                      \tag{2}
\]
No coherent choice of the EPPA extensions is needed:
for each permutation of \(G\), any extension maps the
profile of \(w\) to its prescribed translate. Where the
profiles are all distinct, this even determines the
induced action on \(B\) independently of the extension.


Let \(p\) be the number of coordinates
of \(P(w)\) whose sizes lie strictly between
\(0\) and \(t\). If \(2\le p<s\),
the orbit size is at least
\(\binom sp t^p\ge st^2>st\).
If \(p=s\ge2\), it is at least
\(t^s>st\), unless \(s=t=2\).
When \(p=1\), say its size is \(k\),
and \(r\) of the other coordinates are
full, the exact orbit size is
\[
        s\binom tk\binom{s-1}{r}.
\]
The bound (2) forces
\(k\in\{1,t-1\}\), and the remaining
coordinates are either all empty or all
full. This orbit has exactly \(st\)
points, so every vertex of \(B\) has
one of these profiles, with each profile
occurring exactly once. The action of \(W\)
on \(B\) is its natural action on \(s\)
blocks of \(t\) points. Its unordered
pair-orbits are same-block and
different-block pairs. Consequently
\(H[B]\) is one of
\(\varnothing,K_{st},sK_t,\overline{sK_t}\).
Equation (1) leaves only \(H[B]=sK_t\).
The cross-edges connect each \(D_i\)
to exactly one \(C_i\) by a matching
(or its complement), and to every other
\(C_j\) uniformly by no edges or all
edges. Thus \(H\) is a generalised
clique cover.

When \(s=t=2\) and \(p=2\), the four
profiles instead choose one vertex from
\emph{each} of \(C_1,C_2\): they are
the four corners of a square.
The \(W=S_2\wr S_2\) action on these
corners has unordered pair-orbits of
lengths four (sides) and two
(opposite corners). Equation (1)
forces \(H[B]\) to be the
opposite-corner matching. Labelling
the corners \((a,b)\in\{0,1\}^2\),
take \(D_1=\{(0,0),(1,1)\}\) and
\(D_2=\{(0,1),(1,0)\}\).
With these orders the connection
matrix is exactly
\[
 A=\begin{pmatrix}
 M&M\\M&\overline M
 \end{pmatrix},
\]
which is the Wagner graph.

Finally suppose that \emph{every}
outside profile has \(p=0\).
Then every \(C_i\) is contained
in a true-twin class of \(H\).
These classes all have common size
\(a\ge t\), since \(H\) is vertex-transitive.
The quotient graph \(Q\) on twin
classes is vertex-transitive,
has \(2st/a\le2s\) vertices,
and contains the \(s\) classes
meeting \(G\) as an independent set.
Since \(H\) is not homogeneous,
\(Q\) is not edgeless. A
positive-degree regular graph has no
independent set larger than half its
order. Thus \(a=t\) and
\(Q\) has exactly \(2s\) vertices.
The first \(s\) twin classes comprise
\(G\), and the others comprise \(B\).
Equation (1) implies no edges between
distinct latter classes, so
\(H[B]=sK_t\).

Each \(D_j\), being one such twin
class, is connected \emph{completely}
to precisely \(r\) of the \(C_i\).
Regularity makes \(r\) independent
of \(j\). Permutations of the
\(C_i\)'s extend to automorphisms
of \(H\) and permute the \(D_j\)'s,
so all \(\binom sr\) possible
\(r\)-subsets of \([s]\) must be
represented among the \(s\) cliques
\(D_j\). Thus \(\binom sr\le s\),
giving \(r\in\{0,1,s-1,s\}\).
For \(r=0,s\), every cross-entry is
constant. For \(r=1,s-1\), the
\(s\) profiles are precisely all
singletons or co-singletons, each
once. Pair the \(C_i,D_i\)
using their distinguished coordinate;
all diagonal entries are equal to
\(0\) or \(1\), and so are all
off-diagonal entries. This again is a
generalised clique cover. \(\square\)

\paragraph{Why this is useful.}
Together with the earlier complementary-half
profile lemma, the proof certifies the
\emph{conclusion} of Lemma 3.7 in
the non-homogeneous setting without using
any of Claims 3.8--3.11.
It is optional: the short individual repair
to Claim 3.9 and the Wagner paragraph can be
kept if the authors prefer minimal changes.
The proof also strengthens the conclusion
by dispensing with imprimitive and complementary
partition hypotheses.

\paragraph{The homogeneous case of Lemma 3.7.}
By Gardiner's classification, a homogeneous graph of even order is
a disjoint union of equal cliques or its complement (the two sporadic
homogeneous graphs have odd orders five and nine).
If \(H=mK_a\), each of the \(s\) selected \(K_t\)'s occupies
a different host component. Since \(B=sK_t\), a host component
contains either only outside vertices (forcing \(a=t\)), or
exactly \(t\) selected and \(t\) outside vertices
(forcing \(a=2t\)). These give \(H=2sK_t\) or
\(H=sK_{2t}\), both uniform generalised covers.
If \(H=\overline{mK_a}\), embedding \(sK_t\) requires
\(s=1\) or \(t=1\): any two vertices from different
selected components must lie in one multipartite part,
which cannot simultaneously contain an edge of \(G\).
For \(s=1\), both selected and outside halves are cliques,
so every multipartite part has size at most two;
thus \(H=K_{2t}\) or \(\overline{tK_2}\).
For \(t=1\), both halves are independent,
so \(H\) is edgeless or \(K_{s,s}\).
All these are uniform generalised clique covers.
This finishes the original Lemma 3.7 at paper level,
under its stated partition hypothesis, including homogeneous hosts.
Only the phrase ``or its complement'' next to the Wagner
matrix needs correction: the displayed graph is 3-regular
whereas its graph complement is 4-regular.
\paragraph{Validation status.}
All short-profile orbit types are already
tested by \texttt{checks/section3\_profiles.py}.
A further independent exhaustive checker of
the exceptional \(s=t=2\) case is in
\texttt{checks/cover\_2k2\_all.py}.
The theorem above is a fully spelled-out
paper proof; its group-action and true-twin
reduction are \emph{not yet proved in Lean}.
No manuscript source has been edited.
\end{document}
